%% ───────────────────────────────────────────────────────────────────
\chapter{Il Sistema Italiano: Dalle Origini al NDC}
%% [SINTESI STORICA -- ~12-15 pagine come da indicazione del relatore]
%% ───────────────────────────────────────────────────────────────────

%% ─── PROMPT LLM ─────────────────────────────────────────────────
%% Scrivi il Capitolo 2 di una tesi magistrale in economia pubblica.
%% ATTENZIONE: il relatore ha indicato che questa parte deve essere
%% SINTETICA (~12-15 pagine). Non è un capitolo storico, è un capitolo
%% di CONTESTO. Ogni sottosezione deve essere funzionale a spiegare
%% il PERCHÉ delle criticità attuali, non a descrivere eventi storici.
%% Tono: denso, essenziale, zero verbosità.
%% Fonti principali: MEF/RGS, Fornero (2018), Lavoce.info,
%%                  Rosen-Gayer-Rapallini cap. 12.
%% ────────────────────────────────────────────────────────────────

\section{Dalle origini alla crisi del sistema retributivo}

    %% PROMPT: Spiega l'evoluzione storica in modo funzionale al
    %% problema attuale. Segui questo filo:
    %% - Perché si è scelto il PAYG nel dopoguerra? (necessità di
    %%   pagare subito le pensioni senza accumulare riserve)
    %% - Il metodo retributivo (DB, 1970-1995): cos'è, perché era
    %%   generoso (pensione = 2% × anni × retribuzione finale),
    %%   perché ha creato insostenibilità latente (le promesse erano
    %%   finanziate dal debito implicito, non dai contributi).
    %% - La crisi: mostra la correlazione storica riforme-debito/PIL
    %%   (grafico Favero, Lavoce.info 2018): ogni sistema generoso
    %%   si è associato ad aumento del debito pubblico.
    %% - Spiega la pressione demografica usando la formula:
    %%   P_{t+1} = c*S_t*(1+m)*(1+n) → quando n ↓, le pensioni
    %%   diventano insostenibili senza ridurre P o aumentare c.
    \subsection{La scelta del PAYG nel dopoguerra:
                ragioni politiche ed economiche}
    \subsection{Il metodo retributivo (DB, 1970--1995):
                meccanismo, generosità e insostenibilità latente}
    \subsection{La pressione demografica e fiscale:
                l'invecchiamento nell'equazione $P = cS(1+m)(1+n)$}
    \subsection{La correlazione storica riforme--debito pubblico:
                1960--2018 come argomento per l'urgenza di riformare}

\section{La stagione delle riforme strutturali (1992--2011)}

    %% PROMPT: Tratta le tre riforme principali in modo comparato.
    %% Per ciascuna: contesto, misura principale, logica economica,
    %% impatto su sostenibilità e adeguatezza.
    %%
    %% RIFORMA AMATO (1992): primo intervento di emergenza.
    %%   - Modifica al calcolo retributivo (da 5 a 10 anni per i privati)
    %%   - Innalzamento età pensionabile (da 60 a 65 anni uomini)
    %%   - Contesto: crisi valutaria, pressione UE
    %%
    %% RIFORMA DINI (1995): la svolta strutturale.
    %%   - Introduce il NDC per i nuovi entrati dal 1/1/1996
    %%   - Il regime misto pro-rata per chi aveva < 18 anni di anzianità
    %%   - La logica: legare la pensione ai contributi versati
    %%     e alla crescita del PIL, non alle ultime retribuzioni
    %%   - Problema: effetti a regime solo dopo 30-40 anni
    %%
    %% RIFORMA FORNERO (2011): completamento e urgenza.
    %%   - NDC per tutti per gli anni residui (anche il regime misto)
    %%   - Aggancio automatico dell'età pensionabile alla speranza di vita
    %%   - Contesto: crisi spread 2011, pressione BCE/FMI
    %%   - Paradosso: la riforma "salva i conti" ma crea il problema
    %%     delle pensioni future inadeguate per i giovani
    %%
    %% Costruisci una tabella comparativa delle tre riforme:
    %% | Riforma | Anno | Misura chiave | Impatto sostenibilità |
    %% |         |      |               | Impatto adeguatezza   |
    \subsection{La Riforma Amato (1992):
                primo intervento di emergenza sul retributivo}
    \subsection{La Riforma Dini (1995):
                introduzione del NDC e logica del montante contributivo}
    \subsection{La Riforma Fornero (2011):
                NDC per tutti e aggancio automatico alla speranza di vita}
    \subsection{Il regime misto ``a tre velocità'':
                retributivo puro, misto, contributivo puro --
                tre generazioni, tre destini pensionistici diversi}

\section{L'architettura attuale: meccanismi e anomalie}

    %% PROMPT: Descrivi il sistema come funziona OGGI (2026).
    %% Tre blocchi:
    %%
    %% BLOCCO 1 -- Meccanismo NDC in dettaglio:
    %%   - Aliquota contributiva: 33% dipendenti, 20% autonomi
    %%   - Interesse implicito g = media mobile quinquennale PIL nominale
    %%   - Coefficiente di trasformazione: valori 2024 (4,270% a 67 anni
    %%     fino a 6,655% a 71 anni), aggiornamento biennale
    %%   - Come si calcola concretamente la pensione (formula)
    %%
    %% BLOCCO 2 -- Le deroghe politiche:
    %%   - Quota 100 (2019-2021): 62 anni + 38 contributivi
    %%   - Quota 102 (2022), Quota 103 (2023-in corso)
    %%   - Perché queste deroghe violano la logica NDC: vanno in pensione
    %%     prima che il montante sia sufficiente, riducendo la pensione
    %%     propria E generando costo per il sistema
    %%   - Costo stimato di Quota 100: ~8 miliardi/anno (INPS)
    %%
    %% BLOCCO 3 -- I numeri della spesa:
    %%   - 15,4% PIL nel 2024, 37,2% di tutte le uscite correnti (UPB 2025)
    %%   - La "gobba" MEF nel 2044 (picco per pensionamento Baby Boomers)
    %%   - La discesa strutturale post-2044 per entrata a regime del NDC
    %%   - Gli scenari demografici ISTAT integrati nelle proiezioni
    \subsection{Come funziona il calcolo NDC oggi:
                aliquote, coefficienti 2024 e aggiornamento biennale}
    \subsection{Le deroghe politiche alla logica NDC:
                Quota 100, 102, 103 tra populismo e rigore attuariale}
    \subsection{La spesa pensionistica al 15,4\% del PIL (2024),
                il picco del 2044 e la discesa strutturale post-2044}
    \subsection{Il paradosso redistributivo:
                le pensioni riducono il Gini di $-10{,}7$ punti,
                più di IRPEF e trasferimenti in natura insieme
                (Banca d'Italia, 2025)}

